\chapter{MÔ PHỎNG, KẾT QUẢ VÀ THẢO LUẬN}
\label{chap:ket-qua}

\section{Mục tiêu và nguyên tắc đánh giá}
\label{sec:muc-tieu-danh-gia}

Chương này đánh giá sáu phương pháp trên cùng mô hình vật lý: TD3, DDPG, PPO, AO--SCA,
AO--Grid và AnalyticalRIS. Thay vì tìm một phương pháp thắng duy nhất, kết quả được đọc
theo hai nhóm chính. Nhóm DRL gồm hai thuật toán học ngoài chính sách là DDPG và TD3, cùng
một thuật toán học theo chính sách hiện hành là PPO. Nhóm tối ưu lặp gồm AO--SCA và
AO--Grid. AnalyticalRIS được giữ như một đối chiếu giải tích nhẹ, có chi phí tính toán rất
thấp nhưng không tối ưu đầy đủ.

Ba câu hỏi đánh giá được đặt ra. Thứ nhất, khi dùng cùng biểu diễn hành động có cấu trúc,
các thuật toán DRL đạt mức tổng tốc độ và QoS như thế nào, và có khác biệt ổn định giữa
hai nhóm học ngoài chính sách và học theo chính sách hiện hành hay không? Thứ hai, chất
lượng nghiệm của nhóm DRL cách nhóm AO bao xa? Thứ ba, đổi lại cho chênh lệch chất lượng
đó, độ trễ ra quyết định giảm được bao nhiêu khi dùng một lượt truyền xuôi của mạng thay
cho tối ưu lặp.

Mọi so sánh chất lượng dùng cùng $1.000$ kịch bản kiểm thử tại từng $N$. Đối với DRL, báo
cáo chính là trung bình của năm hạt giống huấn luyện. Với AO và AnalyticalRIS, phương pháp
là xác định nên không có biến thiên do huấn luyện; độ biến thiên của chúng đến từ kịch bản
kênh. Hai loại bất định này được giữ tách biệt trong phần diễn giải thống kê.

\section{Thiết lập mô phỏng}
\label{sec:thiet-lap-mo-phong}

Bảng~\ref{tab:ch4-thiet-lap} tóm tắt cấu hình đã mô tả ở Chương~\ref{chap:phuong-phap}.
Các giá trị được giữ cố định cho tất cả phương pháp khi có ý nghĩa tương ứng.

\begin{table}[htbp]
\centering
\caption{Thiết lập mô phỏng chính}
\label{tab:ch4-thiet-lap}
\footnotesize
\begin{tabularx}{\textwidth}{@{}X l@{}}
\toprule
\textbf{Thành phần} & \textbf{Giá trị} \\
\midrule
Số UE & $4$ \\
Số phần tử STAR--RIS & $16,\ 32,\ 64,\ 96,\ 128$ \\
Công suất phát chuẩn hóa & $1$ \\
Công suất tạp âm & $10^{-3}$ \\
QoS tối thiểu mỗi UE & $0,5$ bit/s/Hz \\
Hệ số tỷ lệ kênh trực tiếp, BS--RIS và RIS--UE & $0,35 / 1,0 / 0,75$ \\
Ngân hàng huấn luyện, xác thực và kiểm thử & $10.000 / 1.000 / 1.000$ kịch bản \\
Số hạt giống huấn luyện cho mỗi thuật toán DRL & $5$ \\
Ngân sách mỗi hạt giống & $100.000$ tương tác môi trường \\
Chu kỳ xác thực & $5.000$ bước \\
Số mẫu đánh giá kiểm thử & $1.000$ kịch bản trên mỗi $N$ \\
\bottomrule
\end{tabularx}
\tablesource{Tác giả tổng hợp từ cấu hình và ngân hàng kịch bản của thí nghiệm (2026)}
\end{table}

Một điểm quan trọng của giao thức là các kịch bản kiểm thử được ghép cặp giữa các phương
pháp. Nghĩa là nếu một kịch bản cụ thể có kênh thuận lợi hoặc khó, tất cả phương pháp đều
được đánh giá trên đúng kịch bản đó. Cách làm này giúp phép so sánh theo kịch bản loại bớt
một phần biến thiên do dữ liệu.

\section{Diễn biến trên tập xác thực trong quá trình huấn luyện}
\label{sec:duong-cong-xac-thuc}

Trước khi so sánh chất lượng cuối, Hình~\ref{fig:ch4-duong-cong} cho thấy diễn biến tổng
tốc độ trên tập xác thực theo số bước huấn luyện. Mỗi đường là trung bình của năm hạt
giống, còn dải tô nhạt là khoảng giữa giá trị nhỏ nhất và lớn nhất trong năm hạt giống đó.

Ba khung cho thấy ba dáng học khác nhau. TD3 tăng rất nhanh trong khoảng $5.000$ bước đầu
rồi gần như đi ngang, với dải hạt giống hẹp ở mọi kích thước. DDPG cũng lên nhanh nhưng
xuất hiện vài lần sụt giảm rồi hồi phục ở giai đoạn giữa và cuối, dấu hiệu quen thuộc của
việc Critic đánh giá quá cao trong một số lần cập nhật. PPO đi lên chậm và đều hơn, cần
gần như toàn bộ ngân sách $100.000$ bước để tiến tới vùng bão hòa, và có dải hạt giống
rộng nhất trong giai đoạn đầu.

\begin{figure}[htbp]
\centering
\resizebox{\textwidth}{!}{\begin{tikzpicture}[font=\small]
\begin{axis}[name=td3, width=6.0cm, height=6.2cm, title={TD3}, title style={font=\small}, xlabel={Bước huấn luyện}, ylabel={Tổng tốc độ trên tập xác thực (bit/s/Hz)}, xmin=0, xmax=100000, ymin=3, ymax=23, scaled x ticks=false, xtick={0,50000,100000}, xticklabels={$0$,$50\text{k}$,$100\text{k}$}, grid=both, grid style={gray!25}, tick label style={font=\scriptsize}, label style={font=\scriptsize}]
\addplot[cUEc!25,draw=none,forget plot] coordinates {(0,4.3243) (5000,13.0227) (10000,13.3760) (15000,13.5739) (20000,13.5755) (25000,13.7876) (30000,13.7883) (35000,13.9466) (40000,14.0038) (45000,14.0164) (50000,14.2563) (55000,14.2341) (60000,14.2964) (65000,14.2854) (70000,14.2975) (75000,14.3467) (80000,14.3226) (85000,14.3057) (90000,14.3411) (95000,14.3308) (100000,14.3700) (100000,14.6551) (95000,14.6814) (90000,14.6181) (85000,14.6673) (80000,14.6110) (75000,14.6670) (70000,14.6217) (65000,14.5735) (60000,14.6257) (55000,14.6068) (50000,14.6907) (45000,14.5468) (40000,14.6045) (35000,14.5132) (30000,14.4455) (25000,14.5261) (20000,14.4001) (15000,14.4136) (10000,14.4193) (5000,14.3190) (0,4.3777)} \closedcycle;
\addplot[very thick,cUEc] coordinates {(0,4.3513) (5000,13.8501) (10000,14.0785) (15000,14.1924) (20000,14.1770) (25000,14.2536) (30000,14.2586) (35000,14.3192) (40000,14.3594) (45000,14.3922) (50000,14.4816) (55000,14.4754) (60000,14.4789) (65000,14.4654) (70000,14.4770) (75000,14.4942) (80000,14.4846) (85000,14.4875) (90000,14.4869) (95000,14.5024) (100000,14.5208)};
\addplot[cUEa!25,draw=none,forget plot] coordinates {(0,4.3117) (5000,15.3049) (10000,15.2975) (15000,15.4240) (20000,15.4026) (25000,15.4155) (30000,15.4850) (35000,15.5742) (40000,15.6513) (45000,15.6819) (50000,15.6614) (55000,15.6804) (60000,15.6675) (65000,15.7366) (70000,15.7182) (75000,15.7594) (80000,15.7970) (85000,15.8021) (90000,15.7805) (95000,15.0497) (100000,15.7798) (100000,16.4020) (95000,16.4909) (90000,16.5224) (85000,16.3488) (80000,16.4149) (75000,16.3503) (70000,16.3497) (65000,16.3788) (60000,16.3164) (55000,16.2530) (50000,16.2462) (45000,16.1716) (40000,16.1815) (35000,16.1005) (30000,16.1325) (25000,16.0279) (20000,15.9923) (15000,15.9673) (10000,15.9829) (5000,15.7456) (0,4.3747)} \closedcycle;
\addplot[very thick,cUEa] coordinates {(0,4.3504) (5000,15.5934) (10000,15.7141) (15000,15.7892) (20000,15.8051) (25000,15.8455) (30000,15.8967) (35000,15.9550) (40000,15.9973) (45000,15.9962) (50000,16.0422) (55000,16.0696) (60000,16.0843) (65000,16.1154) (70000,16.1214) (75000,16.1730) (80000,16.1851) (85000,16.1981) (90000,16.1831) (95000,15.9659) (100000,16.1890)};
\addplot[cUEb!25,draw=none,forget plot] coordinates {(0,4.3349) (5000,17.6507) (10000,17.6426) (15000,17.7530) (20000,17.8450) (25000,17.8486) (30000,17.9048) (35000,18.0327) (40000,18.0402) (45000,18.0947) (50000,18.1936) (55000,18.1760) (60000,18.0985) (65000,18.3406) (70000,16.9783) (75000,18.3046) (80000,18.5115) (85000,18.2800) (90000,18.4317) (95000,18.2653) (100000,18.5449) (100000,18.9046) (95000,18.8666) (90000,18.9639) (85000,18.7874) (80000,18.7197) (75000,18.8972) (70000,18.6179) (65000,18.7468) (60000,18.5448) (55000,18.6235) (50000,18.6336) (45000,18.5043) (40000,18.5697) (35000,18.4368) (30000,18.4672) (25000,18.3582) (20000,18.3481) (15000,18.2146) (10000,18.1099) (5000,17.9844) (0,4.3657)} \closedcycle;
\addplot[very thick,cUEb] coordinates {(0,4.3541) (5000,17.8351) (10000,17.9498) (15000,18.0144) (20000,18.0678) (25000,18.0696) (30000,18.1434) (35000,18.1970) (40000,18.2426) (45000,18.3008) (50000,18.3514) (55000,18.3657) (60000,18.3064) (65000,18.4830) (70000,18.1733) (75000,18.5913) (80000,18.6251) (85000,18.6267) (90000,18.7179) (95000,18.6097) (100000,18.7310)};
\addplot[cUEd!25,draw=none,forget plot] coordinates {(0,4.3255) (5000,18.6185) (10000,18.7576) (15000,18.9503) (20000,18.9222) (25000,18.9589) (30000,19.0903) (35000,19.1965) (40000,19.1605) (45000,19.2716) (50000,19.2850) (55000,19.3061) (60000,19.3812) (65000,19.5341) (70000,19.5273) (75000,19.5819) (80000,19.6335) (85000,19.6988) (90000,19.6911) (95000,17.6183) (100000,19.5794) (100000,20.0405) (95000,19.9375) (90000,19.9722) (85000,19.9311) (80000,19.9536) (75000,19.9783) (70000,19.8592) (65000,19.8567) (60000,19.8045) (55000,19.5847) (50000,19.6557) (45000,19.5463) (40000,19.5695) (35000,19.4783) (30000,19.4798) (25000,19.4265) (20000,19.4584) (15000,19.3656) (10000,19.2871) (5000,19.1409) (0,4.3684)} \closedcycle;
\addplot[very thick,cUEd] coordinates {(0,4.3455) (5000,18.9124) (10000,18.9643) (15000,19.0654) (20000,19.1643) (25000,19.2027) (30000,19.2646) (35000,19.3162) (40000,19.3612) (45000,19.4059) (50000,19.4404) (55000,19.4654) (60000,19.5701) (65000,19.6463) (70000,19.7059) (75000,19.7359) (80000,19.7724) (85000,19.8217) (90000,19.8297) (95000,19.3958) (100000,19.8524)};
\addplot[cChung!25,draw=none,forget plot] coordinates {(0,4.3181) (5000,19.1797) (10000,18.9930) (15000,19.5421) (20000,19.6501) (25000,19.8873) (30000,19.6382) (35000,20.0335) (40000,20.2181) (45000,20.2239) (50000,20.2858) (55000,20.3727) (60000,20.4155) (65000,20.4650) (70000,20.5555) (75000,20.5798) (80000,20.6761) (85000,20.6043) (90000,20.8263) (95000,8.7672) (100000,8.6530) (100000,20.9528) (95000,20.9712) (90000,20.9181) (85000,20.9692) (80000,20.9109) (75000,20.9357) (70000,20.7874) (65000,20.8206) (60000,20.8435) (55000,20.8012) (50000,20.7665) (45000,20.6864) (40000,20.5851) (35000,20.4633) (30000,20.3939) (25000,20.3812) (20000,20.3308) (15000,20.3158) (10000,20.2111) (5000,20.1860) (0,4.3897)} \closedcycle;
\addplot[very thick,cChung] coordinates {(0,4.3571) (5000,19.8352) (10000,19.8426) (15000,19.9930) (20000,20.0479) (25000,20.1285) (30000,20.1299) (35000,20.3367) (40000,20.4540) (45000,20.5190) (50000,20.5532) (55000,20.5904) (60000,20.6535) (65000,20.6580) (70000,20.7180) (75000,20.7756) (80000,20.8188) (85000,20.7957) (90000,20.8735) (95000,18.4484) (100000,18.4220)};
\end{axis}
\begin{axis}[name=ddpg, width=6.0cm, height=6.2cm, title={DDPG}, title style={font=\small}, xlabel={Bước huấn luyện}, yticklabel=\empty, xmin=0, xmax=100000, ymin=3, ymax=23, scaled x ticks=false, xtick={0,50000,100000}, xticklabels={$0$,$50\text{k}$,$100\text{k}$}, grid=both, grid style={gray!25}, tick label style={font=\scriptsize}, label style={font=\scriptsize}, at={(td3.east)}, anchor=west, xshift=0.8cm, legend style={font=\small,at={(0.5,-0.26)},anchor=north,draw=gray!50,legend columns=5,/tikz/every even column/.append style={column sep=10pt}}, legend cell align=left]
\addplot[cUEc!25,draw=none,forget plot] coordinates {(0,4.3243) (5000,13.5532) (10000,13.8025) (15000,13.6749) (20000,13.7603) (25000,13.8853) (30000,14.1248) (35000,14.0132) (40000,14.0553) (45000,14.1519) (50000,14.1167) (55000,14.1509) (60000,14.1502) (65000,14.1291) (70000,14.1612) (75000,14.1911) (80000,14.2250) (85000,14.2319) (90000,14.2708) (95000,14.3137) (100000,14.2758) (100000,14.8759) (95000,14.8956) (90000,14.8979) (85000,14.8902) (80000,14.8972) (75000,14.8832) (70000,14.8604) (65000,14.7898) (60000,14.8294) (55000,14.8132) (50000,14.8687) (45000,14.7434) (40000,14.8027) (35000,14.8122) (30000,14.7274) (25000,14.7473) (20000,14.6081) (15000,14.6267) (10000,14.5812) (5000,14.5536) (0,4.3777)} \closedcycle;
\addplot[very thick,cUEc] coordinates {(0,4.3513) (5000,14.1336) (10000,14.2830) (15000,14.2647) (20000,14.2896) (25000,14.3865) (30000,14.4835) (35000,14.4804) (40000,14.4764) (45000,14.5038) (50000,14.5430) (55000,14.5456) (60000,14.5435) (65000,14.4723) (70000,14.5504) (75000,14.5837) (80000,14.5774) (85000,14.5815) (90000,14.5957) (95000,14.6450) (100000,14.6288)};
\addlegendentry{$N=16$}
\addplot[cUEa!25,draw=none,forget plot] coordinates {(0,4.3117) (5000,15.7533) (10000,15.9218) (15000,16.0619) (20000,16.2019) (25000,16.2224) (30000,16.2323) (35000,16.3299) (40000,12.7116) (45000,16.4230) (50000,16.3773) (55000,16.5003) (60000,8.2600) (65000,16.5032) (70000,16.1640) (75000,16.4890) (80000,16.2163) (85000,16.2731) (90000,16.4987) (95000,16.6603) (100000,16.6621) (100000,16.9233) (95000,16.9430) (90000,16.9365) (85000,16.9153) (80000,16.8904) (75000,16.7997) (70000,16.8907) (65000,16.7974) (60000,16.7192) (55000,16.7804) (50000,16.7842) (45000,16.7684) (40000,16.5953) (35000,16.7019) (30000,16.5891) (25000,16.5817) (20000,16.5511) (15000,16.5037) (10000,16.4847) (5000,16.4066) (0,4.3747)} \closedcycle;
\addplot[very thick,cUEa] coordinates {(0,4.3504) (5000,16.1198) (10000,16.2048) (15000,16.2860) (20000,16.3786) (25000,16.4054) (30000,16.4488) (35000,16.5052) (40000,15.7424) (45000,16.5561) (50000,16.5818) (55000,16.6158) (60000,14.2589) (65000,16.6907) (70000,16.6417) (75000,16.6974) (80000,16.6501) (85000,16.7268) (90000,16.8089) (95000,16.8373) (100000,16.8524)};
\addlegendentry{$N=32$}
\addplot[cUEb!25,draw=none,forget plot] coordinates {(0,4.3349) (5000,17.7315) (10000,17.8802) (15000,17.9043) (20000,18.0110) (25000,18.1823) (30000,18.2700) (35000,18.3833) (40000,18.4872) (45000,18.3015) (50000,18.4376) (55000,18.5138) (60000,18.5818) (65000,15.0607) (70000,13.9785) (75000,18.7781) (80000,18.7850) (85000,18.7560) (90000,18.8170) (95000,18.8548) (100000,18.0683) (100000,19.0698) (95000,19.0430) (90000,19.0604) (85000,19.0295) (80000,19.0716) (75000,18.9580) (70000,18.8800) (65000,18.9845) (60000,18.8898) (55000,18.9327) (50000,18.8390) (45000,18.6987) (40000,18.7529) (35000,18.5996) (30000,18.6225) (25000,18.5563) (20000,18.4662) (15000,18.4192) (10000,18.4420) (5000,18.2461) (0,4.3657)} \closedcycle;
\addplot[very thick,cUEb] coordinates {(0,4.3541) (5000,17.9728) (10000,18.1306) (15000,18.2469) (20000,18.3339) (25000,18.4177) (30000,18.4704) (35000,18.4984) (40000,18.5842) (45000,18.5552) (50000,18.6250) (55000,18.7024) (60000,18.7187) (65000,17.8212) (70000,17.6109) (75000,18.8637) (80000,18.9149) (85000,18.9135) (90000,18.9448) (95000,18.9642) (100000,18.7753)};
\addlegendentry{$N=64$}
\addplot[cUEd!25,draw=none,forget plot] coordinates {(0,4.3255) (5000,18.7184) (10000,19.0179) (15000,19.0762) (20000,19.1920) (25000,19.2957) (30000,19.3375) (35000,19.4353) (40000,19.4810) (45000,19.5890) (50000,19.5831) (55000,19.5849) (60000,19.6451) (65000,19.6761) (70000,8.4095) (75000,19.8101) (80000,19.7694) (85000,19.8781) (90000,19.7775) (95000,18.7665) (100000,19.8574) (100000,20.3482) (95000,20.3010) (90000,20.1974) (85000,20.2159) (80000,20.2184) (75000,20.1684) (70000,20.0297) (65000,20.0898) (60000,19.9965) (55000,19.8621) (50000,19.8197) (45000,19.7723) (40000,19.7426) (35000,19.6346) (30000,19.6090) (25000,19.5937) (20000,19.4292) (15000,19.4490) (10000,19.2781) (5000,19.2530) (0,4.3684)} \closedcycle;
\addplot[very thick,cUEd] coordinates {(0,4.3455) (5000,19.0684) (10000,19.1818) (15000,19.2844) (20000,19.3341) (25000,19.4461) (30000,19.4590) (35000,19.5332) (40000,19.6080) (45000,19.6615) (50000,19.6868) (55000,19.7083) (60000,19.7960) (65000,19.8724) (70000,17.5594) (75000,19.9307) (80000,19.9221) (85000,19.9929) (90000,19.9610) (95000,19.8274) (100000,20.0716)};
\addlegendentry{$N=96$}
\addplot[cChung!25,draw=none,forget plot] coordinates {(0,4.3181) (5000,19.7071) (10000,19.8873) (15000,19.9831) (20000,20.0473) (25000,20.1073) (30000,20.1597) (35000,20.2636) (40000,20.3162) (45000,20.5065) (50000,20.5113) (55000,9.6623) (60000,20.6458) (65000,19.4029) (70000,18.1516) (75000,20.7458) (80000,20.8011) (85000,8.6080) (90000,20.7498) (95000,20.8567) (100000,20.8973) (100000,21.1075) (95000,21.0644) (90000,21.0827) (85000,21.0205) (80000,20.9852) (75000,21.0612) (70000,21.0278) (65000,20.8519) (60000,20.9230) (55000,20.9306) (50000,20.8491) (45000,20.7254) (40000,20.7714) (35000,20.7572) (30000,20.5764) (25000,20.5871) (20000,20.5885) (15000,20.5453) (10000,20.5401) (5000,20.3605) (0,4.3897)} \closedcycle;
\addplot[very thick,cChung] coordinates {(0,4.3571) (5000,19.9977) (10000,20.1446) (15000,20.2143) (20000,20.2875) (25000,20.3210) (30000,20.4113) (35000,20.4990) (40000,20.5327) (45000,20.6100) (50000,20.6403) (55000,18.5175) (60000,20.7533) (65000,20.4998) (70000,20.2559) (75000,20.8725) (80000,20.8894) (85000,18.4139) (90000,20.9157) (95000,20.9570) (100000,20.9697)};
\addlegendentry{$N=128$}
\end{axis}
\begin{axis}[name=ppo, width=6.0cm, height=6.2cm, title={PPO}, title style={font=\small}, xlabel={Bước huấn luyện}, yticklabel=\empty, xmin=0, xmax=100000, ymin=3, ymax=23, scaled x ticks=false, xtick={0,50000,100000}, xticklabels={$0$,$50\text{k}$,$100\text{k}$}, grid=both, grid style={gray!25}, tick label style={font=\scriptsize}, label style={font=\scriptsize}, at={(ddpg.east)}, anchor=west, xshift=0.8cm]
\addplot[cUEc!25,draw=none,forget plot] coordinates {(0,5.9402) (6144,5.7737) (10240,5.7718) (16384,5.7754) (20480,5.7431) (26624,5.7179) (30720,5.7281) (36864,5.8273) (40960,5.8915) (45056,5.9991) (51200,6.1314) (55296,6.1654) (61440,6.4276) (65536,5.9828) (71680,6.0490) (75776,6.0720) (81920,6.0886) (86016,6.1134) (90112,6.1095) (96256,6.1709) (100000,6.1863) (100000,14.6290) (96256,14.4294) (90112,14.1677) (86016,13.8159) (81920,13.4769) (75776,12.8101) (71680,12.6896) (65536,12.4384) (61440,12.3023) (55296,12.6641) (51200,12.4268) (45056,11.9165) (40960,11.6856) (36864,9.6026) (30720,8.5594) (26624,8.4787) (20480,8.3499) (16384,8.2886) (10240,8.3165) (6144,8.0723) (0,7.0458)} \closedcycle;
\addplot[very thick,cUEc] coordinates {(0,6.6176) (6144,7.0143) (10240,7.0831) (16384,7.1370) (20480,7.1779) (26624,7.2645) (30720,7.3376) (36864,7.6500) (40960,8.7147) (45056,9.0227) (51200,9.3606) (55296,9.4322) (61440,8.9022) (65536,8.9226) (71680,9.0587) (75776,9.0927) (81920,9.4215) (86016,9.6476) (90112,9.8074) (96256,9.9459) (100000,9.9575)};
\addplot[cUEa!25,draw=none,forget plot] coordinates {(0,5.9613) (6144,6.6089) (10240,6.9265) (16384,7.1360) (20480,7.3475) (26624,7.6581) (30720,7.8411) (36864,7.9051) (40960,8.3531) (45056,8.1955) (51200,8.8079) (55296,9.2201) (61440,9.7586) (65536,9.7591) (71680,9.8138) (75776,9.7448) (81920,9.9705) (86016,10.0087) (90112,10.2613) (96256,10.9212) (100000,11.5337) (100000,16.7902) (96256,16.7764) (90112,16.7630) (86016,16.7759) (81920,16.7754) (75776,16.7429) (71680,16.8813) (65536,16.8918) (61440,16.8890) (55296,16.8731) (51200,16.8607) (45056,16.8214) (40960,16.7107) (36864,15.1908) (30720,10.8226) (26624,10.5388) (20480,8.8260) (16384,8.4072) (10240,8.2241) (6144,8.0965) (0,7.2402)} \closedcycle;
\addplot[very thick,cUEa] coordinates {(0,6.6848) (6144,7.3048) (10240,7.5359) (16384,7.7261) (20480,7.9763) (26624,8.6033) (30720,9.0973) (36864,9.8633) (40960,11.7919) (45056,12.1497) (51200,12.9327) (55296,13.2204) (61440,13.6872) (65536,13.8910) (71680,14.2626) (75776,14.0766) (81920,14.4137) (86016,14.5844) (90112,14.8607) (96256,15.1972) (100000,15.3808)};
\addplot[cUEb!25,draw=none,forget plot] coordinates {(0,6.5833) (6144,7.3988) (10240,8.5162) (16384,9.6833) (20480,10.3733) (26624,11.4708) (30720,13.9397) (36864,15.1967) (40960,15.0790) (45056,15.0549) (51200,15.4756) (55296,18.3550) (61440,12.9020) (65536,13.8371) (71680,14.2222) (75776,14.7391) (81920,15.5408) (86016,15.7179) (90112,14.9217) (96256,15.8810) (100000,16.5231) (100000,19.2404) (96256,19.2354) (90112,19.2374) (86016,19.2287) (81920,19.2209) (75776,19.2134) (71680,19.2071) (65536,19.2117) (61440,19.2175) (55296,19.1924) (51200,19.1657) (45056,19.1232) (40960,19.0696) (36864,18.9994) (30720,18.5130) (26624,18.0730) (20480,17.2401) (16384,14.8901) (10240,12.7364) (6144,11.1982) (0,7.9795)} \closedcycle;
\addplot[very thick,cUEb] coordinates {(0,7.3927) (6144,9.8425) (10240,10.9874) (16384,12.6049) (20480,14.0229) (26624,15.7008) (30720,16.6006) (36864,17.6241) (40960,17.9608) (45056,18.1037) (51200,18.2635) (55296,18.8680) (61440,17.6651) (65536,17.8949) (71680,17.9957) (75776,17.7071) (81920,17.9124) (86016,17.9772) (90112,17.3093) (96256,17.6336) (100000,17.8392)};
\addplot[cUEd!25,draw=none,forget plot] coordinates {(0,5.6227) (6144,6.0424) (10240,6.0994) (16384,6.6339) (20480,6.8015) (26624,7.1936) (30720,8.5512) (36864,8.3852) (40960,9.0694) (45056,9.7027) (51200,11.5592) (55296,13.6473) (61440,15.1479) (65536,16.3867) (71680,18.0079) (75776,17.8054) (81920,19.0296) (86016,19.3439) (90112,19.4458) (96256,19.5928) (100000,19.6631) (100000,20.4031) (96256,20.3981) (90112,20.4005) (86016,20.3895) (81920,20.3828) (75776,20.3745) (71680,20.3729) (65536,20.3598) (61440,20.3317) (55296,20.2827) (51200,20.2726) (45056,20.2738) (40960,20.2300) (36864,20.1198) (30720,20.0117) (26624,19.8759) (20480,19.6106) (16384,18.7189) (10240,15.7867) (6144,11.3729) (0,7.4731)} \closedcycle;
\addplot[very thick,cUEd] coordinates {(0,6.3652) (6144,8.3018) (10240,9.9062) (16384,11.4267) (20480,12.0768) (26624,13.3157) (30720,14.2928) (36864,15.3840) (40960,15.8555) (45056,16.2694) (51200,16.9306) (55296,17.6256) (61440,18.3199) (65536,18.7612) (71680,19.3256) (75776,19.4093) (81920,19.8116) (86016,19.9699) (90112,20.0245) (96256,20.0827) (100000,20.1134)};
\addplot[cChung!25,draw=none,forget plot] coordinates {(0,6.1554) (6144,6.8846) (10240,7.2709) (16384,7.7954) (20480,7.8736) (26624,8.7359) (30720,8.5289) (36864,16.2652) (40960,16.6491) (45056,17.8410) (51200,19.1700) (55296,19.6248) (61440,19.7361) (65536,19.9617) (71680,20.4328) (75776,20.4274) (81920,20.5196) (86016,20.5914) (90112,20.5909) (96256,20.7098) (100000,20.7179) (100000,21.1550) (96256,21.1522) (90112,21.1515) (86016,21.1505) (81920,21.1502) (75776,21.1452) (71680,21.1470) (65536,21.1494) (61440,21.1151) (55296,21.0895) (51200,21.1283) (45056,21.1362) (40960,21.0803) (36864,21.0822) (30720,20.9917) (26624,20.8926) (20480,20.4332) (16384,19.2345) (10240,16.9361) (6144,12.0906) (0,7.7542)} \closedcycle;
\addplot[very thick,cChung] coordinates {(0,6.7192) (6144,9.1059) (10240,10.9208) (16384,12.4693) (20480,14.6334) (26624,16.8419) (30720,17.4823) (36864,19.5500) (40960,19.8708) (45056,20.2095) (51200,20.5427) (55296,20.6876) (61440,20.7723) (65536,20.8174) (71680,20.8780) (75776,20.8872) (81920,20.9106) (86016,20.9281) (90112,20.9316) (96256,20.9573) (100000,20.9644)};
\end{axis}
\end{tikzpicture}
}
\caption{Diễn biến tổng tốc độ trên tập xác thực trong quá trình huấn luyện}
\label{fig:ch4-duong-cong}
\figuresource{Tác giả vẽ từ dữ liệu xác thực định kỳ của năm hạt giống (2026)}
\end{figure}

Điểm chung là cả ba thuật toán đều đạt vùng ổn định trước khi hết ngân sách huấn luyện, vì
vậy kết quả kiểm thử ở các mục sau không phải là ảnh chụp của một quá trình còn dang dở.

\section{Độ ổn định của nhóm DRL qua các hạt giống huấn luyện}
\label{sec:do-on-dinh}

Trước khi so sánh với AO, cần xem bản thân các phương pháp DRL ổn định đến đâu giữa năm
hạt giống. Bảng~\ref{tab:ch4-do-lech} trình bày độ lệch chuẩn của tổng tốc độ theo trung
bình của từng hạt giống. Giá trị nhỏ cho thấy kết quả ít phụ thuộc khởi tạo ngẫu nhiên;
giá trị lớn cho thấy các lần huấn luyện có thể hội tụ về chất lượng rất khác nhau.

\begin{table}[htbp]
\centering
\caption{Độ lệch chuẩn tổng tốc độ giữa năm hạt giống của nhóm DRL}
\label{tab:ch4-do-lech}
\footnotesize
\begin{tabular}{@{}lrrr@{}}
\toprule
$N$ & \textbf{TD3} & \textbf{DDPG} & \textbf{PPO} \\
\midrule
$16$ & 0,1713 & 0,3028 & 3,6307 \\
$32$ & 0,2931 & 0,1166 & 2,3044 \\
$64$ & 0,1265 & 0,0939 & 0,2111 \\
$96$ & 0,1134 & 0,2162 & 0,2974 \\
$128$ & 0,0755 & 0,0748 & 0,1716 \\
\bottomrule
\end{tabular}
\tablesource{Tác giả tổng hợp từ năm hạt giống huấn luyện (2026)}
\end{table}


Kết quả cho thấy PPO biến thiên lớn tại $N = 16$ và $N = 32$, sau đó ổn định rõ khi $N \ge
64$. TD3 và DDPG nhìn chung có độ lệch chuẩn thấp hơn trong hầu hết cấu hình. Tuy nhiên,
từ $N = 64$ trở lên, cả ba phương pháp đều có độ biến thiên tương đối nhỏ so với mức tổng
tốc độ khoảng $19$ đến $21$ bit/s/Hz. Điều này cho thấy khi biểu diễn hành động có cấu
trúc tạo được một miền học thuận lợi, khác biệt giữa các thuật toán giảm đáng kể ở kích
thước STAR--RIS lớn hơn.

Không nên diễn giải Bảng~\ref{tab:ch4-do-lech} thành kết luận rằng nhóm học ngoài chính
sách luôn ổn định hơn nhóm học theo chính sách hiện hành. PPO ở $N = 64$ đến $128$ đã có
độ lệch chuẩn cùng bậc với TD3 và DDPG. Kết luận phù hợp hơn là PPO nhạy hơn ở các cấu
hình nhỏ trong bộ thí nghiệm này, còn tại $N$ lớn không có sự chênh lệch ổn định đủ lớn để
khẳng định một họ thuật toán luôn ưu thế.

\section{Hiệu năng tổng tốc độ trên tập kiểm thử}
\label{sec:hieu-nang-tong-toc-do}

Bảng~\ref{tab:ch4-tong-toc-do} trình bày tổng tốc độ trung bình của sáu phương pháp. Đây
là bảng chính để đánh giá chất lượng nghiệm.

\begin{table}[htbp]
\centering
\caption{Tổng tốc độ trung bình trên 1.000 kịch bản kiểm thử (bit/s/Hz)}
\label{tab:ch4-tong-toc-do}
\footnotesize
\begin{tabular}{@{}lrrrrrr@{}}
\toprule
$N$ & \textbf{TD3} & \textbf{DDPG} & \textbf{PPO} & \textbf{AO--SCA} & \textbf{AO--Grid} & \textbf{AnalyticalRIS} \\
\midrule
$16$ & 14,5515 & 14,4912 & 9,9633 & 14,3449 & 16,5802 & 1,9812 \\
$32$ & 16,3264 & 16,8518 & 15,4663 & 17,7256 & 18,5044 & 1,9820 \\
$64$ & 18,7324 & 18,9625 & 19,0146 & 19,9050 & 20,4401 & 1,9821 \\
$96$ & 19,9279 & 20,0913 & 20,1327 & 20,9103 & 21,5862 & 1,9821 \\
$128$ & 20,9195 & 20,9890 & 21,0201 & 21,6095 & 22,3918 & 1,9821 \\
\bottomrule
\end{tabular}
\tablesource{Tác giả tổng hợp từ kết quả kiểm thử đã khóa (2026)}
\end{table}


Hình~\ref{fig:ch4-tong-toc-do} thể hiện cùng dữ liệu dưới dạng đồ thị. Có ba xu hướng dễ
nhận thấy. Thứ nhất, tổng tốc độ của hầu hết phương pháp tối ưu và học đều tăng theo $N$,
phản ánh lợi ích của việc có nhiều phần tử STAR--RIS hơn để tạo kênh hiệu dụng mạnh. Thứ
hai, AO--Grid đạt tổng tốc độ cao nhất tại cả năm kích thước. Thứ ba, ba thuật toán DRL
tiến rất gần nhau ở $N \ge 64$ và không có một thuật toán luôn cao nhất.

\begin{figure}[htbp]
\centering
\resizebox{0.95\textwidth}{!}{\begin{tikzpicture}[font=\small]
\begin{axis}[
  width=13cm, height=8cm,
  xlabel={Số phần tử STAR--RIS $N$},
  ylabel={Tổng tốc độ trung bình (bit/s/Hz)},
  xtick={16,32,64,96,128}, xmin=8, xmax=136,
  ymin=0, ymax=24,
  grid=both, grid style={gray!25},
  tick label style={font=\small}, label style={font=\small},
  legend style={font=\small,at={(0.02,0.98)},anchor=north west,draw=gray!50,
                legend columns=2,/tikz/every even column/.append style={column sep=6pt}},
  legend cell align=left,
]
\addplot[very thick,cUEc,mark=*,mark size=2pt] coordinates {(16,14.551465) (32,16.326370) (64,18.732389) (96,19.927889) (128,20.919524)};
\addlegendentry{TD3}
\addplot[very thick,cUEa,mark=square*,mark size=2pt] coordinates {(16,14.491241) (32,16.851819) (64,18.962516) (96,20.091340) (128,20.989047)};
\addlegendentry{DDPG}
\addplot[very thick,cUEb,mark=triangle*,mark size=2pt] coordinates {(16,9.963296) (32,15.466326) (64,19.014617) (96,20.132726) (128,21.020078)};
\addlegendentry{PPO}
\addplot[very thick,cUEd,mark=diamond*,mark size=2pt] coordinates {(16,14.344949) (32,17.725581) (64,19.904969) (96,20.910255) (128,21.609453)};
\addlegendentry{AO--SCA}
\addplot[very thick,cChung,mark=pentagon*,mark size=2pt] coordinates {(16,16.580238) (32,18.504410) (64,20.440070) (96,21.586214) (128,22.391796)};
\addlegendentry{AO--Grid}
\addplot[very thick,cKenh,mark=o,mark size=2pt] coordinates {(16,1.981170) (32,1.981977) (64,1.982059) (96,1.982070) (128,1.982074)};
\addlegendentry{AnalyticalRIS}
\end{axis}
\end{tikzpicture}
}
\caption{Tổng tốc độ của sáu phương pháp theo số phần tử STAR--RIS}
\label{fig:ch4-tong-toc-do}
\figuresource{Tác giả vẽ từ dữ liệu kiểm thử (2026)}
\end{figure}

AnalyticalRIS gần như không thay đổi theo $N$ và chỉ đạt khoảng $1,98$ bit/s/Hz. Kết quả
này cho thấy quy tắc căn pha đơn giản dùng làm khởi tạo không đủ để khai thác đồng thời
công suất, phân bổ tốc độ luồng chung và QoS. Vì vậy, AnalyticalRIS phù hợp như một đối
chiếu cực nhẹ hơn là một bộ giải chất lượng cao.

\section{So sánh trong nội bộ nhóm DRL}
\label{sec:so-sanh-drl}

Tại $N = 16$, TD3 và DDPG đạt khoảng $14,5$ bit/s/Hz trong khi PPO chỉ khoảng $10$
bit/s/Hz. Tại $N = 32$, DDPG cao nhất trong nhóm DRL. Từ $N = 64$ trở lên, PPO và DDPG
nhỉnh hơn TD3 một lượng nhỏ nhưng khoảng cách chỉ vài phần trăm bit/s/Hz.

Các phép so sánh theo $1.000$ kịch bản có thể cho giá trị $p$ rất nhỏ vì mỗi kịch bản được
ghép cặp và số mẫu lớn. Tuy nhiên, một tuyên bố về thuật toán phải tính đến biến thiên
huấn luyện chỉ có năm hạt giống. Ví dụ tại $N = 32$, chênh lệch giữa TD3 và DDPG là khoảng
$-0,525$ bit/s/Hz theo trung bình; phép kiểm định theo kịch bản cho kết quả rất rõ, nhưng
kiểm định theo năm hạt giống không còn mạnh sau hiệu chỉnh nhiều so sánh. Vì vậy, luận văn
không kết luận DDPG tốt hơn TD3 một cách tổng quát từ một kích thước duy nhất.

Quan sát đáng chú ý hơn là cả ba thuật toán DRL cùng đạt vùng $20$ đến $21$ bit/s/Hz ở
$N = 96$ và $N = 128$. Điều này phù hợp với giả thuyết rằng bộ giải mã hành động đã mã hóa
một phần lớn cấu trúc khó của bài toán. Khi miền hành động trở nên thuận lợi hơn, khác
biệt giữa hai cách cập nhật chính sách không còn chi phối tuyệt đối chất lượng cuối.

\section{So sánh nhóm DRL với nhóm tối ưu lặp}
\label{sec:so-sanh-drl-ao}

AO--Grid cao hơn tất cả các phương pháp DRL ở mọi $N$. AO--SCA cũng cao hơn DRL tại $N \ge
32$, trong khi ở $N = 16$ thì TD3 và DDPG có trung bình nhỉnh hơn AO--SCA một lượng nhỏ.

Kết quả này phù hợp với bản chất hai nhóm phương pháp. AO nhận một CSI và dành nhiều phép
đánh giá hàm mục tiêu để điều chỉnh trực tiếp biến vật lý; nó còn tìm kiếm trên miền rộng
hơn bộ giải mã có cấu trúc của DRL, vì vậy có khả năng đi tới các nghiệm mà Actor không
biểu diễn được. Ngược lại, DRL dành chi phí lớn ở giai đoạn huấn luyện nhưng tại thời điểm
sử dụng chỉ thực hiện một lượt mạng nơ-ron và bộ giải mã.

Để tránh hiểu nhầm, AO--SCA và AO--Grid trong luận văn không được xem là cận trên hay
nghiệm tối ưu toàn cục. Chúng là các bộ giải cục bộ mạnh trong giao thức hiện tại. Việc
AO--Grid cao hơn AO--SCA cũng không mâu thuẫn với lý thuyết: bài toán phi lồi có nhiều
miền cục bộ, và tìm kiếm lưới với các bước lớn, khả năng tắt luồng bằng giá trị $0$ cùng
thứ tự quét thuận nghịch có thể tìm được vùng tốt hơn gradient lân cận.

\section{Đánh giá chất lượng dịch vụ}
\label{sec:danh-gia-qos}

Bảng~\ref{tab:ch4-qos} trình bày xác suất tất cả bốn UE cùng đạt yêu cầu $R_k \ge 0,5$
bit/s/Hz. Đây là chỉ số khắt khe hơn tỷ lệ UE đạt QoS vì chỉ cần một UE vi phạm thì kịch
bản được tính là không đạt toàn bộ.

\begin{table}[htbp]
\centering
\caption{Xác suất tất cả UE đạt QoS}
\label{tab:ch4-qos}
\footnotesize
\begin{tabular}{@{}lrrrrrr@{}}
\toprule
$N$ & \textbf{TD3} & \textbf{DDPG} & \textbf{PPO} & \textbf{AO--SCA} & \textbf{AO--Grid} & \textbf{AnalyticalRIS} \\
\midrule
$16$ & 0,9720 & 0,9848 & 0,7788 & 0,9440 & 1,0000 & 0,0000 \\
$32$ & 0,9898 & 0,9930 & 0,9916 & 0,9980 & 1,0000 & 0,0000 \\
$64$ & 0,9994 & 0,9982 & 0,9986 & 1,0000 & 1,0000 & 0,0000 \\
$96$ & 0,9996 & 0,9992 & 0,9998 & 1,0000 & 1,0000 & 0,0000 \\
$128$ & 0,9996 & 1,0000 & 0,9994 & 1,0000 & 1,0000 & 0,0000 \\
\bottomrule
\end{tabular}
\tablesource{Tác giả tổng hợp từ dữ liệu kiểm thử (2026)}
\end{table}


Hình~\ref{fig:ch4-qos} cho thấy rõ sự hội tụ về vùng QoS cao khi $N$ tăng. Tại $N = 16$,
PPO là trường hợp yếu nhất với khoảng $77,9\%$ kịch bản đạt đồng thời QoS, trong khi TD3,
DDPG và AO--SCA đều trên $94\%$. Từ $N = 32$, hầu hết phương pháp đều vượt $98,9\%$.
AO--Grid đạt $100\%$ trong toàn bộ bộ kết quả báo cáo. AnalyticalRIS không đạt yêu cầu
đồng thời cho cả bốn UE trong bất kỳ kịch bản tổng hợp nào, nên được lược khỏi hình để
giữ thang trục dễ đọc; điều này nhấn mạnh rằng căn pha đơn thuần không đủ nếu không phối
hợp với phân bổ RSMA.

\begin{figure}[htbp]
\centering
\resizebox{0.95\textwidth}{!}{\begin{tikzpicture}[font=\small]
\begin{axis}[
  width=13cm, height=7.5cm,
  xlabel={Số phần tử STAR--RIS $N$},
  ylabel={Xác suất tất cả UE đạt QoS},
  xtick={16,32,64,96,128}, xmin=8, xmax=136,
  ymin=0.74, ymax=1.02,
  grid=both, grid style={gray!25},
  tick label style={font=\small}, label style={font=\small},
  legend style={font=\small,at={(0.5,-0.18)},anchor=north,draw=gray!50,
                legend columns=5,/tikz/every even column/.append style={column sep=10pt}},
  legend cell align=left,
]
\addplot[very thick,cUEc,mark=*,mark size=2pt] coordinates {(16,0.972000) (32,0.989800) (64,0.999400) (96,0.999600) (128,0.999600)};
\addlegendentry{TD3}
\addplot[very thick,cUEa,mark=square*,mark size=2pt] coordinates {(16,0.984800) (32,0.993000) (64,0.998200) (96,0.999200) (128,1.000000)};
\addlegendentry{DDPG}
\addplot[very thick,cUEb,mark=triangle*,mark size=2pt] coordinates {(16,0.778800) (32,0.991600) (64,0.998600) (96,0.999800) (128,0.999400)};
\addlegendentry{PPO}
\addplot[very thick,cUEd,mark=diamond*,mark size=2pt] coordinates {(16,0.944000) (32,0.998000) (64,1.000000) (96,1.000000) (128,1.000000)};
\addlegendentry{AO--SCA}
\addplot[very thick,cChung,mark=pentagon*,mark size=2pt] coordinates {(16,1.000000) (32,1.000000) (64,1.000000) (96,1.000000) (128,1.000000)};
\addlegendentry{AO--Grid}
\end{axis}
\end{tikzpicture}
}
\caption{Xác suất tất cả UE đạt QoS của các phương pháp có chất lượng cạnh tranh}
\label{fig:ch4-qos}
\figuresource{Tác giả vẽ từ dữ liệu kiểm thử (2026)}
\end{figure}

Kết quả QoS cũng giúp giải thích vì sao chỉ nhìn tổng tốc độ là chưa đủ. Một phương pháp
có thể đạt tổng tốc độ cao bằng cách dồn tài nguyên cho UE mạnh, nhưng nếu xác suất đạt
QoS đồng thời thấp thì không phù hợp mục tiêu. Bộ chọn điểm lưu mô hình của DRL ưu tiên
tính khả thi QoS trước khi tối đa hóa tổng tốc độ nhằm hạn chế hiện tượng này.

\section{Độ trễ ra quyết định}
\label{sec:do-tre}

Bảng~\ref{tab:ch4-do-tre} trình bày độ trễ tính toán trên cùng một bộ xử lý trung tâm
(Central Processing Unit, CPU) một luồng. Giá trị báo cáo là trung vị trên $100$ lần đo cho
mỗi cấu hình. Đây là thời gian của chính sách hoặc bộ giải, không phải độ trễ đầu cuối của
hệ thống vô tuyến.

\begin{table}[htbp]
\centering
\caption{Độ trễ ra quyết định trên bộ xử lý trung tâm một luồng, giá trị trung vị (ms)}
\label{tab:ch4-do-tre}
\footnotesize
\begin{tabular}{@{}lrrrrrr@{}}
\toprule
$N$ & \textbf{TD3} & \textbf{DDPG} & \textbf{PPO} & \textbf{AO--SCA} & \textbf{AO--Grid} & \textbf{AnalyticalRIS} \\
\midrule
$16$ & 0,3389 & 0,3435 & 0,5645 & 830,4545 & 188,5626 & 0,2031 \\
$32$ & 0,3660 & 0,3652 & 0,5847 & 1064,7201 & 320,0277 & 0,2094 \\
$64$ & 0,3934 & 0,3930 & 0,6237 & 1861,6429 & 579,2715 & 0,2209 \\
$96$ & 0,4343 & 0,4328 & 0,6850 & 2543,5247 & 881,0458 & 0,2230 \\
$128$ & 0,4519 & 0,4459 & 0,7252 & 3606,3791 & 1204,7310 & 0,2349 \\
\bottomrule
\end{tabular}
\tablesource{Tác giả tổng hợp từ phép đo trên bộ xử lý trung tâm một luồng (2026)}
\end{table}


\begin{figure}[htbp]
\centering
\resizebox{0.95\textwidth}{!}{\begin{tikzpicture}[font=\small]
\begin{semilogyaxis}[
  width=13cm, height=8cm,
  xlabel={Số phần tử STAR--RIS $N$},
  ylabel={Độ trễ ra quyết định (ms, thang log)},
  xtick={16,32,64,96,128}, xmin=8, xmax=136,
  ymin=0.1, ymax=8000,
  grid=both, grid style={gray!25},
  tick label style={font=\small}, label style={font=\small},
  legend style={font=\small,at={(0.5,-0.18)},anchor=north,draw=gray!50,
                legend columns=3,/tikz/every even column/.append style={column sep=10pt}},
  legend cell align=left,
]
\addplot[very thick,cUEc,mark=*,mark size=2pt] coordinates {(16,0.338859) (32,0.365976) (64,0.393359) (96,0.434294) (128,0.451860)};
\addlegendentry{TD3}
\addplot[very thick,cUEa,mark=square*,mark size=2pt] coordinates {(16,0.343531) (32,0.365179) (64,0.393010) (96,0.432756) (128,0.445910)};
\addlegendentry{DDPG}
\addplot[very thick,cUEb,mark=triangle*,mark size=2pt] coordinates {(16,0.564507) (32,0.584740) (64,0.623697) (96,0.684953) (128,0.725154)};
\addlegendentry{PPO}
\addplot[very thick,cUEd,mark=diamond*,mark size=2pt] coordinates {(16,830.454527) (32,1064.720090) (64,1861.642865) (96,2543.524655) (128,3606.379085)};
\addlegendentry{AO--SCA}
\addplot[very thick,cChung,mark=pentagon*,mark size=2pt] coordinates {(16,188.562568) (32,320.027747) (64,579.271550) (96,881.045757) (128,1204.731037)};
\addlegendentry{AO--Grid}
\addplot[very thick,cKenh,mark=o,mark size=2pt] coordinates {(16,0.203087) (32,0.209429) (64,0.220871) (96,0.222986) (128,0.234883)};
\addlegendentry{AnalyticalRIS}
\draw[gray!60,dashed,thick] (axis cs:8,1) -- (axis cs:136,1);
\node[font=\scriptsize,gray!45!black,anchor=south east] at (axis cs:134,1.1) {ngưỡng $1$ ms};
\end{semilogyaxis}
\end{tikzpicture}
}
\caption{Độ trễ ra quyết định khi số phần tử STAR--RIS tăng}
\label{fig:ch4-do-tre}
\figuresource{Tác giả vẽ từ phép đo trên bộ xử lý trung tâm một luồng (2026)}
\end{figure}

Ba thuật toán DRL đều duy trì độ trễ dưới $1$ ms ở toàn bộ miền $N$. DDPG và TD3 gần như
tương đương vì dùng cùng kiến trúc Actor xác định. PPO chậm hơn một chút do xử lý chính
sách Gaussian nhưng vẫn ở mức dưới một mili giây. AO--Grid cần từ khoảng $189$ ms đến hơn
$1,2$ giây; AO--SCA từ khoảng $0,83$ giây đến $3,61$ giây. AnalyticalRIS nhanh nhất vì chỉ
thực hiện công thức giải tích và các ánh xạ đơn giản.

Kết luận hợp lý không phải là DRL nhanh nhất, bởi AnalyticalRIS còn nhanh hơn. Kết luận là
DRL tạo được vùng đánh đổi hấp dẫn: chất lượng tổng tốc độ gần AO hơn rất nhiều so với
AnalyticalRIS, trong khi độ trễ lại gần AnalyticalRIS hơn nhiều so với AO.

\section{Đánh đổi giữa chất lượng và độ trễ}
\label{sec:danh-doi}

Hình~\ref{fig:ch4-danh-doi} đặt tổng tốc độ và độ trễ cạnh nhau tại $N = 128$. Trục ngang
dùng thang logarit để thể hiện khoảng cách vài bậc độ lớn giữa DRL và AO.

\begin{figure}[htbp]
\centering
\resizebox{0.95\textwidth}{!}{\begin{tikzpicture}[font=\small]
\begin{semilogxaxis}[
  width=13cm, height=8cm,
  xlabel={Độ trễ ra quyết định (ms, thang log)},
  ylabel={Tổng tốc độ (bit/s/Hz)},
  xmin=0.1, xmax=20000, ymin=0, ymax=26,
  grid=both, grid style={gray!25},
  minor grid style={gray!12},
  tick label style={font=\small}, label style={font=\small},
]
\addplot[only marks,cUEc,mark=*,mark size=3.4pt] coordinates {(0.4519,20.9195)};
\node[font=\small,cUEc,anchor=east,xshift=-6pt,yshift=-9pt,inner sep=1pt] at (axis cs:0.4519,20.9195) {TD3};
\addplot[only marks,cUEa,mark=square*,mark size=3.4pt] coordinates {(0.4459,20.9890)};
\node[font=\small,cUEa,anchor=east,xshift=-6pt,yshift=9pt,inner sep=1pt] at (axis cs:0.4459,20.9890) {DDPG};
\addplot[only marks,cUEb,mark=triangle*,mark size=3.4pt] coordinates {(0.7252,21.0201)};
\node[font=\small,cUEb,anchor=west,xshift=6pt,yshift=0pt,inner sep=1pt] at (axis cs:0.7252,21.0201) {PPO};
\addplot[only marks,cUEd,mark=diamond*,mark size=3.4pt] coordinates {(3606.3791,21.6095)};
\node[font=\small,cUEd,anchor=north,xshift=0pt,yshift=-8pt,inner sep=1pt] at (axis cs:3606.3791,21.6095) {AO--SCA};
\addplot[only marks,cChung,mark=pentagon*,mark size=3.4pt] coordinates {(1204.7310,22.3918)};
\node[font=\small,cChung,anchor=south,xshift=0pt,yshift=8pt,inner sep=1pt] at (axis cs:1204.7310,22.3918) {AO--Grid};
\addplot[only marks,cKenh,mark=o,mark size=3.4pt] coordinates {(0.2349,1.9821)};
\node[font=\small,cKenh,anchor=west,xshift=8pt,yshift=0pt,inner sep=1pt] at (axis cs:0.2349,1.9821) {AnalyticalRIS};
\draw[gray!55,dashed] (axis cs:1,0) -- (axis cs:1,26);
\node[font=\scriptsize,gray!45!black,anchor=south,rotate=90] at (axis cs:0.92,7.5) {ngưỡng $1$ ms};
\node[font=\scriptsize,gray!45!black,anchor=north west] at (axis cs:0.11,25.6) {vùng độ trễ thấp};
\node[font=\scriptsize,gray!45!black,anchor=north east] at (axis cs:18000,18.6) {vùng độ trễ cao};
\draw[gray!50,-{Stealth}] (axis cs:2.0,4.2) -- (axis cs:700,4.2);
\node[font=\scriptsize,gray!45!black,anchor=south] at (axis cs:38,4.3)
     {chậm hơn ba đến bốn bậc độ lớn};
\end{semilogxaxis}
\end{tikzpicture}
}
\caption{Không gian đánh đổi giữa tổng tốc độ và độ trễ tại $N = 128$}
\label{fig:ch4-danh-doi}
\figuresource{Tác giả vẽ từ dữ liệu kiểm thử và phép đo độ trễ (2026)}
\end{figure}

Hình cho thấy không có một phương pháp trội hoàn toàn trên cả hai trục. AO--Grid nằm cao
nhất về tổng tốc độ nhưng xa về bên phải do độ trễ lớn. AnalyticalRIS nằm rất gần bên trái
nhưng thấp về chất lượng. Ba thuật toán DRL tạo thành một cụm ở vùng phía trên bên trái:
tổng tốc độ khoảng $21$ bit/s/Hz với thời gian dưới $1$ ms. Điểm của TD3 và DDPG gần như
trùng nhau vì độ trễ chỉ lệch $0,006$ ms và tổng tốc độ lệch $0,07$ bit/s/Hz. Vì vậy, nếu
ứng dụng có chu kỳ CSI ngắn, DRL có thể là lựa chọn thực tế dù không đạt tổng tốc độ cao
nhất.

\section{Khả năng mở rộng theo số phần tử STAR--RIS}
\label{sec:kha-nang-mo-rong}

Khi $N$ tăng từ $16$ lên $128$, số chiều hành động tăng từ $57$ lên $393$ và số chiều
trạng thái cũng tăng đáng kể. Dù vậy, độ trễ của Actor DRL chỉ tăng khoảng vài phần mười
mili giây vì phép tính chủ yếu là hai lớp kết nối đầy đủ có kích thước cố định $256$, còn
lớp đầu và lớp cuối mới thay đổi theo số chiều. Ngược lại, thời gian của AO tăng mạnh vì
số tọa độ STAR--RIS cần thử hoặc tính gradient tăng, và mỗi vòng lặp cần nhiều lần đánh
giá hàm vật lý.

Tổng tốc độ của nhóm DRL và AO đều tăng theo $N$. Điều này cho thấy việc thêm phần tử
STAR--RIS vẫn mang lại lợi ích trong mô hình. Khoảng cách tuyệt đối giữa AO--Grid và cụm
DRL tại $N = 128$ vào khoảng $1,37$ đến $1,47$ bit/s/Hz, trong khi chênh lệch độ trễ là ba
đến bốn bậc độ lớn. Đây là bằng chứng trực quan nhất cho lập luận về đánh đổi giữa chất
lượng và độ trễ của luận văn.

\section{Phân tích thống kê và phạm vi diễn giải}
\label{sec:phan-tich-thong-ke}

Bộ kết quả lưu cả so sánh ghép cặp theo kịch bản và so sánh theo hạt giống huấn luyện. Hai
đơn vị này trả lời hai câu hỏi khác nhau. Phép ghép cặp theo $1.000$ kịch bản hỏi rằng với
các chính sách đã huấn luyện và cố định, phương pháp A có thường cho kết quả cao hơn B
trên cùng kênh hay không. Phép so sánh theo năm hạt giống hỏi rằng nếu xem quá trình huấn
luyện là nguồn ngẫu nhiên, chênh lệch giữa hai thuật toán có ổn định qua các lần huấn
luyện hay không.

Ví dụ tại $N = 16$, TD3 cao hơn DDPG trung bình khoảng $0,060$ bit/s/Hz. Với $1.000$ kịch
bản ghép cặp, phép kiểm định $t$ cho giá trị nhỏ và có thể xem chênh lệch của hai chính
sách là có hệ thống. Tuy nhiên, khi ghép theo năm hạt giống, giá trị $p$ lớn và không còn
ý nghĩa sau hiệu chỉnh Holm. Vì vậy, luận văn không chuyển kết quả ở mức kịch bản thành
kết luận rằng TD3 tốt hơn DDPG như một thuật toán.

Tại $N = 32$, AO--Grid cao hơn cả TD3 và DDPG khoảng $2,18$ và $1,65$ bit/s/Hz. Chênh lệch
này đủ lớn để tồn tại cả ở phân tích theo kịch bản và theo hạt giống. Đây là trường hợp có
thể kết luận mạnh hơn rằng trong giao thức hiện tại, AO--Grid cho chất lượng tổng tốc độ
cao hơn nhóm DRL ở kích thước đó.

Các chỉ số QoS nằm trong $[0, 1]$, trong khi khoảng tin cậy chuẩn có thể vượt ra ngoài
miền này khi trung bình rất gần $1$. Vì vậy, các khoảng tin cậy trong dữ liệu thô được xem
là mô tả thống kê chứ không phải xác suất vật lý vượt $100\%$. Nếu mở rộng nghiên cứu, phép
lấy mẫu lặp hoặc mô hình nhị thức sẽ phù hợp hơn cho các chỉ số tỷ lệ.

Ngoài ra, số hạt giống bằng năm còn tương đối nhỏ cho các tuyên bố tổng quát về khác biệt
thuật toán. Kết quả đủ để mô tả hành vi trong giao thức luận văn, nhưng chưa đủ để khẳng
định thứ hạng phổ quát của TD3, DDPG và PPO trên các mô hình kênh khác.

\section{Vai trò của biểu diễn hành động có cấu trúc}
\label{sec:vai-tro-bieu-dien}

Một dấu hiệu quan trọng là khi dùng chung bộ giải mã, DDPG, TD3 và PPO tiến rất gần nhau
tại $N$ lớn dù cơ chế cập nhật khác nhau. Điều này ủng hộ nhận định rằng biểu diễn hành
động có vai trò lớn. Bộ giải mã đã làm ba việc khó thay cho mạng: giới hạn công suất luồng
chung vào vùng ưu tiên, biến chênh lệch đầu ra thành phân bổ công suất riêng gần như dồn
hết vào một luồng, và cung cấp pha tham chiếu theo CSI.

Tuy nhiên, kết quả hiện tại chưa phải một phân tích tách thành phần đầy đủ. Các phép thử loại bỏ (ablation) có sẵn gồm không dùng RIS, RIS cố định, RIS ngẫu nhiên và chia công suất đều; chúng
kiểm tra vai trò của RIS và công suất nhưng chưa tách riêng đóng góp của phần công suất có
cấu trúc và của riêng pha tham chiếu. Vì vậy, luận văn trình bày vai trò của biểu diễn
hành động như một kết luận được hỗ trợ bởi phân tích cấu trúc và hành vi thực nghiệm, đồng
thời coi phép tách theo từng thành phần của bộ giải mã là hướng cần bổ sung.

\section{Khi nào nên dùng DRL và khi nào nên dùng AO}
\label{sec:khi-nao-dung}

Nếu CSI thay đổi chậm và hệ thống có đủ thời gian tính toán, AO--Grid hoặc AO--SCA có thể
phù hợp vì đạt tổng tốc độ cao hơn trong thí nghiệm. Nếu CSI thay đổi nhanh hoặc bộ điều
khiển phải ra quyết định dưới một mili giây, DRL có ưu thế rõ sau khi đã trả chi phí huấn
luyện ngoại tuyến (offline). Trong các thiết bị có tài nguyên tính toán rất hạn chế và chỉ cần một
quyết định thô, AnalyticalRIS có thể là lựa chọn nhẹ nhất nhưng phải chấp nhận QoS và tổng
tốc độ thấp trong cấu hình này.

Do đó, luận văn không đề xuất một thuật toán duy nhất cho mọi tình huống. Kết quả cung cấp
một bản đồ lựa chọn: AO ở phía chất lượng cao và độ trễ lớn; AnalyticalRIS ở phía độ trễ
thấp và chất lượng thấp; DRL nằm giữa nhưng nghiêng mạnh về phía độ trễ thấp.

\section{Hạn chế nghiên cứu}
\label{sec:han-che}

Các hạn chế chính gồm mô hình kênh Gaussian tổng hợp, CSI hoàn hảo, SIC lý tưởng, số UE cố
định bằng bốn, năm hạt giống huấn luyện, STAR--RIS cho phép pha truyền và pha phản xạ độc
lập, chưa xét lượng tử hóa pha và phần cứng thực, chưa đo độ trễ đầu cuối, và bộ giải mã
có cấu trúc thu hẹp miền hành động. Các hằng số $0,70$, $0,99$, $10$ và $0,25\pi$ là lựa
chọn thiết kế được cố định trước thực nghiệm, chưa có chứng minh tối ưu lý thuyết.

Một hạn chế khác là hai nhóm DRL và AO không tìm kiếm trên đúng cùng một miền hành động.
AO làm việc trên miền vật lý rộng hơn, còn DRL bị ràng buộc bởi bộ giải mã có cấu trúc.
Đây không phải lỗi của phép so sánh nếu mục tiêu là đánh giá một chính sách triển khai cụ
thể, nhưng phải được nêu rõ để không diễn giải khoảng cách tổng tốc độ như thể chỉ do
thuật toán học kém hơn.

\section{Kết luận chương}
\label{sec:ch4-ket-luan}

Kết quả mô phỏng cho thấy không có bằng chứng rằng TD3 luôn vượt DDPG hoặc PPO. Ở $N$ lớn,
ba thuật toán DRL đạt chất lượng rất gần nhau khi dùng cùng biểu diễn hành động có cấu
trúc. Nhóm AO, đặc biệt là AO--Grid, đạt tổng tốc độ cao hơn nhưng phải trả chi phí tính
toán lớn hơn nhiều. Nhóm DRL duy trì độ trễ dưới một mili giây trong toàn bộ các kích
thước khảo sát và giữ QoS cao từ $N = 32$ trở lên. Vì vậy, đóng góp thực nghiệm chính của
luận văn là làm rõ đánh đổi giữa một chính sách học được có độ trễ thấp và các bộ giải lặp
có chất lượng nghiệm cao hơn.
